% answers_short.tex —— 2014 年·安徽卷·文科 · 简短答案速查
% 由 tools/make_short_answers.py 自 answers.tex 的 \gksolhead/\gkanswer 提取，
% 并与 2014/2014_数学真题_安徽文_答案要点.md 逐题比对后生成，勿手改（改源 MD 后重跑）。

\begin{gknotice}
\textbf{使用说明}：本节只给最终答案便于快速核对，完整推导见解析版。答案为编写者独立推导，\textbf{非官方标准答案}。
\end{gknotice}

\gkpart{选择题}
\begin{center}\begin{tabular}{ccccccccccc}
\toprule
题号 & 1 & 2 & 3 & 4 & 5 & 6 & 7 & 8 & 9 & 10 \\
\midrule
答案 & D & C & A & B & B & D & C & A & D & B \\
\bottomrule
\end{tabular}\end{center}

\gkpart{填空题}
\begin{itemize}[leftmargin=2.6em]
  \item[\textbf{11.}] $\dfrac{27}8$。$\left(\dfrac{16}{81}\right)^{-\frac34}=\left(\dfrac32\right)^3=\dfrac{27}8$，两个对数之和为 $\log_31=0$。
  \item[\textbf{12.}] $\dfrac14$。每一步所得的垂线段与上一条之比恒为 $\dfrac{\sqrt2}2$，$a_7=2\left(\dfrac{\sqrt2}2\right)^6=\dfrac14$。
  \item[\textbf{13.}] $4$。三条边界两两交于 $(0,2)$，$(2,0)$，$(8,-2)$，用行列式计算得面积 $\dfrac12\times8=4$。
  \item[\textbf{14.}] $\dfrac5{16}$。$f\left(\dfrac{29}4\right)=f\left(-\dfrac34\right)=-\dfrac3{16}$，$f\left(\dfrac{41}6\right)=f\left(-\dfrac76\right)=-\sin\dfrac{7\pi}6=\dfrac12$。
  \item[\textbf{15.}] ①③。①中 $x^3$、③中 $\sin x-x$ 在 $P$ 两侧变号；②的 $x=-1$ 不是切线，④中 $\tan x-x$ 两侧同号，⑤中 $\ln x-(x-1)\leqslant0$ 只在单侧。
\end{itemize}

\gkpart{解答题}
\begin{itemize}[leftmargin=2.6em]
  \item[\textbf{16.}] $\cos A=\dfrac13$，$a=2\sqrt2$；或 $\cos A=-\dfrac13$，$a=2\sqrt3$。由 $S=\dfrac12bc\sin A$ 得 $\sin A=\dfrac{2\sqrt2}3$，故 $\cos A=\pm\dfrac13$，再用余弦定理。
  \item[\textbf{17.}] (1) $90$。按女生占比 $\dfrac{4500}{15000}=\dfrac3{10}$ 分层，$300\times\dfrac3{10}=90$；(2) $0.75$。后四组频率之和为 $2\times(0.150+0.125+0.075+0.025)=0.75$；(3) 列联表为男生超过 $4$ 小时 $165$ 人、不超过 $45$ 人，女生分别为 $60$ 人与 $30$ 人；$K^2=\dfrac{100}{21}\approx4.762>3.841$，有 $95\%$ 的把握认为有关
  \item[\textbf{18.}] (1) 见解析版；(2) $S_n=\dfrac{(2n-1)3^{n+1}+3}4$。由（1）得 $a_n=n^2$，$b_n=n\cdot3^n$，用错位相减法求和
  \item[\textbf{19.}] (1) 见解析版；(2) $18$。底面中心到各顶点距离 $4\sqrt2$，由侧棱 $2\sqrt{17}$ 得高 $6$；截面为梯形，$EF=8$，$GH=4$，高 $3$，面积 $\dfrac{8+4}2\times3=18$
  \item[\textbf{20.}] (1) 记 $x_1=\dfrac{-1-\sqrt{4+3a}}3$，$x_2=\dfrac{-1+\sqrt{4+3a}}3$；减区间为 $(-\infty,x_1)$、$(x_2,+\infty)$，增区间为 $(x_1,x_2)$；(2) 最大值点：$0<a<4$ 时为 $x=\dfrac{\sqrt{4+3a}-1}3$，$a\geqslant4$ 时为 $x=1$；最小值点：$0<a<1$ 时为 $x=1$，$a=1$ 时为 $x=0$ 或 $x=1$，$a>1$ 时为 $x=0$
  \item[\textbf{21.}] (1) $5$。由 $|AB|=4$ 与 $|AF_1|=3|BF_1|$ 得 $|AF_1|=3$；周长给出 $|AF_2|+|BF_2|=12$，而两者分别为 $2a-3$ 与 $2a-1$，解得 $a=4$，$|AF_2|=5$；(2) $\dfrac{\sqrt2}2$。记 $|BF_1|=m$，由两个三角形余弦定理与 $|AF_1|=3m$ 得 $\cos\angle AF_1F_2=\dfrac a{2c}$、$m=\dfrac{2(a^2-c^2)}{3a}$，再对 $\triangle ABF_2$ 用余弦定理，得 $4e^4-12e^2+5=0$，取 $e<1$ 得 $e=\dfrac{\sqrt2}2$
\end{itemize}
